Se presentan los resultados obtenidos tras la implementación de la arquitectura de datos y el desarrollo de los tableros de visualización.

\section{Dashboards de Inteligencia de Negocios}

Para satisfacer las necesidades de información financiera y control operativo, se implementan cuatro vistas analíticas principales:

\begin{enumerate}
    \item \textbf{Matriz Jerárquica Interactiva:} Permite la visualización consolidada de la información de gastos desde el nivel general hasta el detalle granular por departamento o cuenta específica.
    \item \textbf{Panel Tabular de Nivel Atómico:} Expone cada transacción individual registrada, integrando folios de movimientos, fechas y referencias cruzadas para auditorías precisas.
    \item \textbf{Vista de Confrontación Contable:} Presenta un resumen totalizado de los montos agrupados por cuenta contable y un desglose por fecha para el cuadre contra la balanza de comprobación.
    \item \textbf{Vista de Validación de Etiquetas Analíticas:} Se utiliza para la validación de las etiquetas analíticas asignadas a cada movimiento, permitiendo que el área financiera identifique las transacciones que carecen de clasificación y establezca las etiquetas correspondientes de manera oportuna.
\end{enumerate}

\begin{figure}[H]
    \centering
    \includegraphics[width=0.8\textwidth]{figuras/resultados/Matriz jerárquica interactiva para consolidación de gastos.png}
    \caption{Matriz jerárquica interactiva para consolidación de gastos.}
    \label{fig:dashboard1}
\end{figure}

\begin{figure}[H]
    \centering
    \includegraphics[width=0.8\textwidth]{figuras/resultados/Panel tabular para validación de transacciones.png}
    \caption{Panel tabular para validación de transacciones.}
    \label{fig:dashboard2}
\end{figure}

\begin{figure}[H]
    \centering
    \includegraphics[width=0.8\textwidth]{figuras/resultados/Vista de Confrontación Contable.png}
    \caption{Vista de Confrontación Contable.}
    \label{fig:dashboard3}
\end{figure}